\documentclass[a4paper,12pt]{article}
\usepackage{graphicx}
\usepackage{hyperref}
\usepackage[italian]{babel}
\usepackage{parskip}
\setlength{\parskip}{1em}
\graphicspath{ {./images/} }
\hypersetup {
    colorlinks=true,
    linkcolor=black,
    pdfborderstyle={/S/U/W 1},
}

\begin{document}
\begin{titlepage}
    \begin{center}
        \begin{minipage}{0.3\textwidth}
            \includegraphics[width=\linewidth]{Logo_Universita_Padova300x300.png}
        \end{minipage}
        
    \end{center}

    \begin{center}
        \noindent\rule{0.7\textwidth}{0.4pt}\\[0.5cm]
        \large\textbf{Università degli Studi di Padova}\\
        \large Corso di Laurea in Informatica\\
        \small A. A. 2025/26
    \end{center}

    \vspace{1,5cm}    

    \begin{center}
        \begin{minipage}{0.5\textwidth}
            \includegraphics[width=\linewidth]{logo.png}
        \end{minipage}
    \end{center}

    \vspace{1,5cm}    
    \begin{center}
        {\Large \textbf{Relazione}\par}
        \vspace{0.6cm}
        {\Large Progetto di Tecnologie Web\par}
        \vspace{0.3cm}
    \end{center}
    \vfill

    \noindent
    \begin{tabular}{r|l}
        \textbf{Studenti} & Sofia De Blasi - matricola - email \\
         & Nenad Radulovic - matricola - email \\
         & Giacomo Speggiorin - matricola - email \\
         & Nicola Simionato - matricola - email \\
         & \\
        \textbf{Indirizzo del sito} & indirizzo\\
         & \\
        \textbf{Utenti} & 
        \textbf{Username:} user, \textbf{Password:} user \\ 
        & \textbf{Username:} admin, \textbf{Password:} admin\\


    \end{tabular}

    \vfill

\end{titlepage}     
\newpage
\tableofcontents
\newpage

\section{Introduzione}
\subsection{Com'è nata l'idea}
L'idea di questo progetto nasce da una domanda che ci siamo posti più volte nella nostra quotidianità: \textit{"Cosa potremmo fare questo weekend?"}. 
Nelle città universitarie vengono spesso organizzati eventi e attività dedicate agli studenti, è però seccante doverne andare ogni volta alla ricerca. 
Bisogna sapere inanzi tutto cosa si sta cercando (il che non è una cosa scontata), e serve poi andarsi a sfogliare di volta in volta le pagine social degli organizzatori, che sono sempre diversi.\par
Abbiamo quindi pensato a quanto sarebbe comodo e utile se fossero gli organizzatori stessi ad andare incontro agli studenti, e a quanto si semplificherebbe la ricerca se esistesse un'unica bacheca dove tutti gli organizzatori possano appendere i loro volantini.\par
Nel mentre buttavamo giù le idee il discorso si è ampliato, e ci siamo resi conto che la fuori ci sono un sacco di servizi di cui gli studenti hanno bisogno ma non hanno un punto di riferimento sul dove andarli a cercare.

\subsection{L'obbiettivo di StudentSpace}
StudentSpace propone una \textit{bacheca online} che raccolga in un unico luogo eventi, attività e servizi di interesse per gli studenti, andando oltre il contesto strettamente accademico.\par
Attraverso la piattaforma sarà possibile trovare eventi sportivi, attività serali, concerti, e altre occasioni di socialità.
Potranno inoltre dare la loro disponibilità ad altri studenti per ripetizioni, oppure consultare annunci di affitto relativi alla propria città.\par
L'obbiettivo principale è garantire a studenti e studentesse un accesso rapido e completo a ciò che potrebbe interessarli.
Allo stesso tempo, si vuole agevolare gli organizzatori, che si rivolgono a questo target, fornendo loro un canale unico e mirato per raggiungere un pubblico ampio e specifico.


\newpage
\section{Analisi del lavoro svolto}
\subsection{Processo decisionale}
Durante tutto lo sviluppo del progetto ci si è trovati a dover prendere diverse decisioni in merito a design, funzionalità, contenuto ecc. 
Nella maggior parte dei casi questi punti suono stati discussi in presenza, tra tutti i membri del gruppo. 
Ovviamente, questo metodo comporta delle chiare limitazioni, si è quindi deciso di adottare
anche dei metodi di comunicazione (sia sincroni che asincroni) alternativi.\\\\
Per quanto riguarda le comunicazioni asincrone, è stata usata ampiamente la piattaforma di messaggistica \textbf{\underline{\nameref{strumenti:whatsapp}}}. 
Si è deciso inoltre di adottare \textbf{\underline{\nameref{strumenti:discord}}} come strumento di comunicazione asincrona, in supporto alle riunioni in presenza.

\subsection{Modalità di lavoro}
Per adattarsi meglio alle esigenze del gruppo si è deciso che i membri del gruppo avrebbero svolto ciascuno le proprie attività in maniera prevalentamente asincrona.
Questa decisione è stata presa per riuscire a dare più libertà ai singoli membri nell'organizzare il proprio lavoro. 
La scelta è stata inoltre supportata dalla volontà del gruppo di dare maggiore fluidità allo sviluppo del progetto. 
Il lavoro è stato infatti suddiviso in modo che ogni membro andasse a lavorare su punti diversi rispetto agli altri, cosicché nessuno dovesse preoccuparsi che il proprio lavoro andasse in conflitto con quello di qualcun altro. 
A sostegno di ciò si è deciso di implementare un sistema di versionamento del codice (con annesso hosting online) tramite \textbf{\underline{\nameref{strumenti:git}}} e \textbf{\underline{\nameref{strumenti:github}}}.

\subsection{Divisione del lavoro}
Visto lo scopo di natura accademica del progetto si è ritenuto opportuno che ciascun membro del gruppo fosse coivolto (quantomeno in parte) in ogni ambito del progetto.
Questo ha permesso a ogni membro di avere un idea generale sul funzionamento del programma, e di mettere mano direttamente a ogni fase del suo sviluppo.


\newpage
\section{Progettazione}
\subsection{Tipi di utente} 
\subsection{Struttura del sito} 
\subsection{Mockup} 
\subsection{Responsive Web Design e Mobile First} 


\newpage
\section{Sviluppo}
\subsection{HTML}
\subsection{CSS}
\subsection{JavaScript}
\subsection{PHP}
\subsection{Database}
\subsection{Errori}
\subsection{Validazione} 
\subsection{Immagini} 
\subsection{Colori} 
\subsection{Font} 
\subsection{.htaccess}

\newpage
\section{Test}
\subsection{Color contrast}

\newpage
\section{Accessibilità}
\subsection{Introduzione} 
\subsection{WAI-ARIA e screen-reader} 
\subsection{Progressive Enhancement} 
\subsection{Supporto ai dispositivi “legacy”} 
\subsection{Ulteriori funzionalità di accessibilità}

\newpage
\section{SEO e prestazioni}

\newpage
\section{Strumenti utilizzati}

\subsection{Canva} \label{strumenti:canva}
\underline{\url{https://www.canva.com/it_it/}} \\\\
Piattaforma di progettazione grafica per creare contenuti visuali per presentazioni, siti web e prodotti simili. 

\subsection{Discord} \label{strumenti:discord}
\underline{\url{https://discord.com/}} \\\\
Piattaforma di messaggistica istantanea, di distribuzione digitale e di voice over IP, utilizzata per la comunicazione all'interno server Discord.
Gli utenti comunicano con chiamate vocali, videochiamate, messaggi di testo, media e file in chat private o come membri di un server Discord. Quest'ultimi sono una raccolta di canali di tipo vocale e/o testuale. 

\subsection{Git} \label{strumenti:git}
\underline{\url{https://git-scm.com/}} \\\\
Applicazione software per il controllo di versione distribuito utilizzabile da interfaccia a riga di comando. Viene utilizzato il versionamento delle repository dedicate al codice e alla documentazione del progetto.

\subsection{Github} \label{strumenti:github}
\underline{\url{https://git-scm.com/}} \\\\
È un servizio di hosting per progetti software che implementa lo strumento di controllo versione distribuito Git.  

\subsection{LaTeX} \label{strumenti:latex}
\underline{\url{https://www.latex-project.org/get/}} \\\\
Linguaggio di marcatura del testo, usato ampiamente in ambito accademico per la preparazione di testi.

\subsection{Silktide} \label{strumenti:silktide}
\underline{\url{https://silktide.com/}} \\\\
È un'estensione che facilita i test sull'accessibilità.

\subsection{W3C Validator} \label{strumenti:w3cvalidator}
\underline{\url{http://validator.w3.org/}} \\\\
Validatore markup HTML/XHTML.

\subsection{WhatsApp} \label{strumenti:whatsapp}
\underline{\url{https://www.whatsapp.com/}} \\\\

\subsection{WebAIM} \label{strumenti:webaim}
\underline{\url{https://webaim.org/resources/contrastchecker/}} \\\\

\newpage
\section{Guida all’installazione del progetto}


\end{document}

